% Appendix C -- physical-model details (moved from Section IV and the limitations table of
% Section VIII-C for length, 2026-10-12; text unchanged apart from joins). Module paths for the
% lead's check are in the comments of sections/04_physical_model.tex (sremu commit 2a016a6).
% Every number is a macro (numbers.yaml; `code` macros are read from the sremu code) or cited.
\section{Physical-Model Details}\label{app:model}

This appendix extends Section~\ref{sec:model}. \claim{C-MODELCARD}Table~\ref{tab:modelcard}
lists every component of the model with its validity domain, verification evidence and
default; the identifiers refer to Section~\ref{sec:vv} and Appendix~\ref{app:vv}.

% GENERATED by sremu/docs/tools/gen_model_card.py from docs/model_card.yaml.
% Do not edit; edit the YAML and rerun the paper pipeline.
\begin{table*}[t]
\caption{Model card of sremu: model, validity domain, verification evidence and default of each component. Generated from docs/model\_card.yaml; evidence ids refer to the test tables of Section V. ENG marks an engineering choice without a literature source.}
\label{tab:modelcard}
\centering
\footnotesize
\setlength{\tabcolsep}{3pt}
\begin{tabular}{@{}>{\raggedright\arraybackslash}p{0.95in}>{\raggedright\arraybackslash}p{1.65in}>{\raggedright\arraybackslash}p{2.25in}>{\raggedright\arraybackslash}p{1.35in}>{\raggedright\arraybackslash}p{0.6in}@{}}
\toprule
Component & Model & Validity domain & Evidence & Default \\
\midrule
LOS & Friis + patterns + explicit V & Far field of the antennas & \texttt{PH-01} & On \\
RIS & K1, \texttt{point} & Cell in the far field of Tx and Rx ($r \geq 3\lambda$); pitch $\leq \lambda/2$; scalar co-pol & \texttt{PH-02--07}, \texttt{G1--3}, \texttt{V-TANG} & auto \\
RIS & K1, \texttt{po\_patch} & As above; uniform-patch PO; Richardson-converged & \texttt{PH-04}, \texttt{PH-26} & Pitch $> \lambda/2$ \\
$\Gamma$ tables & Complex-linear LUT & Inside the grid; band $\leq$ validity bandwidth; reciprocal & \texttt{PH-28}, \texttt{PH-31} & --- \\
Ground & Flat-earth image + Fresnel & Smooth, flat; no curvature (P.526 experimental) & \texttt{PH-10}, \texttt{PH-32}, \texttt{G8} & Off \\
Walls & Image, order 1 & Finite polygons; specular & \texttt{PH-08} & Off \\
Diffraction & Single knife edge $F(\nu)$ & One dominant edge & \texttt{PH-11} & Off \\
Diffuse & Additive SoS & Statistical background only & \texttt{ST-02}, \texttt{G6} & Off \\
Mobility/Doppler & Chord delay + exact phase + tiles & $\varepsilon_{\mathrm{rel}} \leq 10^{-3}$ (bound plus \texttt{verify\_interp}) & \texttt{PH-17--19} & Auto \\
Wideband & FD Kaiser 8/16/32 & $B \leq 0.8 f_s$; $\leq -60$ dB in band & \texttt{PH-20} & Auto \\
Noise & $k T_0 F f_s$ per RX & --- & \texttt{PH-21}, \texttt{G5} & On \\
Non-idealities & Ordered pipeline: command, quantizer, failures, table, efficiency, coupling hook, transition & Tables, failures, efficiency; coupling via tables only & \texttt{PH-13--15}, \texttt{PH-25}, \texttt{PH-30} & --- \\
Transitions & Per-class $g_c$, complex-linear $\Gamma$ path & [ENG] until measured & \texttt{RT-12} & Step \\
\bottomrule
\end{tabular}
\end{table*}


\subsection{Conventions and line of sight}
% contracts.py conventions; physics/los.py, antennas.py, diffraction.py, ground.py, walls.py,
% diffuse.py.
\claim{C-STACK}The trajectories are static, piecewise-linear waypoints, constant velocity, one-
and two-dimensional random waypoint~\cite{Bettstetter2003RWP} with keyed legs, a cubic-spline
trace, or timed updates from a client, so positions never depend on an update grid. Tables
without angle dependence use a separable form of
the coefficient tensor. Tables with angle dependence are cuts in one plane of incidence: a
direction $\mathbf d$ from a cell enters as the signed angle
$\mathrm{atan2}(\mathbf d\cdot\mathbf u,\mathbf d\cdot\mathbf n)$ in the plane of the panel
normal $\mathbf n$ and its axis $\mathbf u$, and its component along the other axis is
projected out (a panel option refuses off-plane geometry instead), while the obliquity and the
front-side indicator of K1 use the full cosines. Tables of the incidence angle alone,
$g(\theta)$, enter through the reciprocal extension $\Gamma(\theta_i,\theta_s)=g(|\theta_i-\theta_s|/2)$, an engineering choice (ENG). A table is read
at the carrier of the link (linear in frequency between its grid points) and held over the
signal's band, which may not exceed the table's validity bandwidth; there is no $\Gamma$ filter.
No measured angle-dependent table is used in this paper. Channel coefficients are dimensionless and normalized as in the Friis
equation, $P_R=P_T|h|^2$, with antenna gains inside $h$. The complex baseband signal $x(t)$
of a carrier $f_c$ is $s(t)=\sqrt2\,\mathrm{Re}\{x(t)e^{j2\pi f_ct}\}$, and a path with delay
$\tau$ maps $x(t)$ to $a\,x(t-\tau)\,e^{-j2\pi f_c\tau}$.

\claim{C-LOS}The LOS path is $h_{\mathrm{LOS}}=\sqrt{G_TG_R}\,
\lambda/(4\pi d)\,V$ with scalar co-polar antenna patterns (isotropic, a $\cos^\alpha$
pattern with peak gain $2(\alpha+1)$, or a tabulated pattern with provenance) and an explicit
visibility factor $V$: zero behind an opaque occluder, $10^{-L/20}$ through a partition with
transmission loss $L$~dB (an engineering choice), or the Fresnel knife-edge factor $F(\nu)$
of one edge, which keeps its phase and adds coherently to the RIS paths. Occluders and
partitions act on the LOS path and on the ground and wall images, never on a RIS hop, so a
blocked scenario whose occluder also crosses RIS hops is idealized (Appendix~\ref{app:cases}).
Scenario names never imply blockage. Optional mechanisms add a flat-ground image path with Fresnel
coefficients, image expansion of the two RIS hops over that ground, order-one images in
finite walls and an additive diffuse term (Table~\ref{tab:modelcard}).

\subsection{Kernel guards, cell models and non-idealities}
% scenario/guards.py CELL_NEAR / CELL_TOO_CLOSE; ris_kernel.energy_budget, KIRCHHOFF_WARNING;
% physics/quadrature.py; physics/nonideal.py.
\claim{C-K1-GUARDS}A scenario check warns when a cell is closer to a node than
three wavelengths, where the point-scatterer approximation weakens, and rejects distances
below one wavelength. With the default obliquity the scattered power of a passive panel does
not exceed the power it intercepts in any tested configuration (PH-23); the alternative obliquity
$(\cos\theta_i+\cos\theta_s)/2$ can exceed it for anomalous reflection and therefore
emits a validity warning when selected. With that Kirchhoff obliquity K1 is the
physical-optics integral of the aperture; with the default it is an area-tied aperture kernel
whose obliquity follows the measured cosine pattern, and both agree at specular reflection.
The continuous physical-optics reference of Fig.~\ref{fig:F3}(a) integrates the same kernel
with the default obliquity.

\claim{C-CELL-MODEL}Forcing \texttt{point} on cells larger than the \texttt{auto} threshold
emits a warning, and tables whose measurement already contains the cell pattern force
\texttt{point} without obliquity, so that the pattern is not counted twice. The two models are
not numerically equivalent: a patch carries the sub-aperture factor of a uniform patch, which a
point omits. \claim{C-RICHARDSON}For \texttt{po\_patch}
panels a Richardson check computes the received power of up to \numRichardsonLinks{}
representative links at sub-cell size $\Delta$ and at $\Delta/2$, halves $\Delta$ at most
\numRichardsonHalvings{} times while the change exceeds \numRichardsonTolDb~dB, and fails
otherwise; the accepted size and the estimate go to the manifest.

\claim{C-CELL-MODEL}The laboratory pitch of \numLabPitchMm~mm lies \numLabPitchMarginMm~mm
(\numLabPitchMarginPct~\%) below half a wavelength at \numLabFreqGHz~GHz, so \texttt{auto}
resolves to \texttt{point} in E-B, in the laboratory-like geometry of E-C and in E-E, whose
cells have the same pitch; the synthetic panel of E-C has a pitch of exactly half a
wavelength, which \texttt{auto} also resolves to \texttt{point}. E-A sets \texttt{point} on its
quarter-wavelength cells, the geometries of Tang et al.\ and E-D resolve to it, and the
laboratory check draws \texttt{point} and \texttt{po\_patch} with equal prior weight. Forcing \texttt{po\_patch} on every panel of E-B and of the E-C laboratory
geometry (coefficient level; record in Appendix~\ref{app:provenance}) lowers the RIS power by
\numCellSensRisMinDb--\numCellSensRisMaxDb~dB. The exact optima of E-B change in at most
\numCellSensCfgPct~\% of the cells, and the \texttt{point} optimum loses at most
\numCellSensDeployLossDb~dB under \texttt{po\_patch}; the change of RX1's channel at the E-B1
switch moves by at most \numCellSensSwitchMaxDb~dB, the most with the line of sight present;
and the E-C losses relative to continuous phase move by at most \numCellSensEcMaxDb~dB. In the
laboratory check the paired M2 predictions of the two models differ by a median of at most
\numLabCmMedianMaxDb~dB per ring (95th percentile at most \numLabCmQninetyfiveMaxDb~dB). For a
single half-wavelength cell at \numCellRetroAngleB$^\circ$ or wider retro-reflection the
default Richardson check does not converge from sub-cells of \numSubcellMax$\lambda$ and fails,
so the value quoted in Section~\ref{sec:model-cells} comes from a finer start.

\claim{C-NONIDEAL}A command (state indices, phases or a codebook entry) is expanded from
controls to physical cells by the panel's control granularity. Stuck and dead cells follow a
failure map, and dead cells absorb by default. The efficiency factor scales the amplitude by
$\sqrt{\eta}$; it is refused for measured or digitized tables unless explicitly overridden,
because their data already contain the loss. The transition of
Section~\ref{sec:design-time} acts in time.

\claim{C-OPTIMUM}For a direction $\theta$ the best per-element choice is
$s_n(\theta)=\arg\max_s\mathrm{Re}\{Z[s,n]e^{-j\theta}\}$; it can change at the critical
angles $\arg(Z[a,n]-Z[b,n])\pm\pi/2$ and nowhere else, and the optimum is the choice at the
direction of the optimal sum, which the sweep over the sorted critical angles finds in
$O(S^2N\log(S^2N))$ operations for $S$ states and $N$ cells.

\subsection{Interpolation bound and fractional delay}
% channel/interp.py (EPS_REL split D1, t_curv_s, tile_side_max_m, choose_policy, M_MAX,
% T_U_MAX_S, VERIFY_INTERP_STRIDE, VERIFY_INTERP_FLAG, verify_interp); channel/fd.py (Kaiser
% sinc, FD_ROWS, N_PHASES polyphase table with linear interpolation, INTEGER_MAX_PHASE_RAD).
\claim{C-DOPPLER}For static scenes snapshots are evaluated at events, and configuration
switches are discontinuities at exact sample indices. Between two snapshots
the delay of a path without a RIS is a chord, linear in time between its delays at the two
snapshots, and its amplitude is complex-linear; the residual of a tile is the sum of its
cells' coefficients relative to the tile's reference delay, interpolated complex-linearly.
\claim{C-INTERP-BOUND}Half of the target $\varepsilon_{\mathrm{rel}}$ goes to each source of
error. The chord error of a path with range $r$ is $k\ddot rT_u^2/8$ with $\ddot r\le
v_\perp^2/r_{\min}$, which gives
\begin{equation}
T_{\mathrm{curv}}=\sqrt{8(\varepsilon_{\mathrm{rel}}/2)\,r_{\min}/(k\,v_\perp^2)}.
\label{eq:tcurv}
\end{equation}
The residual of a tile with side $D_t$ turns by $\Delta\phi$, equal to
$2\pi v_{\mathrm{rel}}D_tT_u/(\lambda R_{\min})$, per interval, and the midpoint dip of linear interpolation,
$\Delta\phi^2/8$, must stay below $\varepsilon_{\mathrm{rel}}/2$, which bounds $D_t$. The
policy starts from $T_{u0}=\min(T_{\mathrm{curv}},T_{u,\max})$ with $T_{u,\max}$ of
\numTuMaxMs~ms (a scenario can lower $T_u$ but not raise it), evaluates the periods
$T_{u0}/2^m$ for $m$ up to \numTuHalvingsMax{} with the tile counts each needs, and picks
the pair with the lowest estimated compute cost; the choice goes to the manifest.
\claim{C-VERIFY-INTERP}Each run checks the bound: in every \numVerifyStride{}th snapshot
interval the engine evaluates the exact snapshot at the interval midpoint and logs
$\max|h_{\mathrm{interp}}-h_{\mathrm{true}}|/|h|$ per link; the manifest keeps the maximum,
and runs above $\numVerifyFlagFactor\,\varepsilon_{\mathrm{rel}}$ are flagged (PH-18).
Section~\ref{sec:vv} compares the rendered Doppler shifts of LOS and RIS paths with closed
forms (PH-17).

\claim{C-FD}Each TX stream declares its occupied bandwidth $B$. A link is narrowband when
$\pi B\Delta\tau\le\numNarrowbandRad$~rad over its delay spread $\Delta\tau$; its paths then
share one envelope delay and keep their own exact carrier phases. Otherwise the paths are
grouped so that each group passes the same test, and each group gets its own
FD filter at its chord delay. The FD filters are Kaiser-windowed sinc
filters ($\beta=\numKaiserBeta$)~\cite{Laakso1996FracDelay} with \numFdTaps{} taps for
$B/f_s$ up to \numFdBands{}, respectively, applied through a polyphase table of
\numFdPhases{} fractional phases with linear interpolation between neighbouring phases.
Their in-band error stays below \numFdTargetDb~dB for every fractional delay (PH-20); a
wider band is rejected. An integer delay without a filter is used when the band is narrow
enough for the half-sample error, $\pi B/(2f_s)$, to stay below \numIntegerDelayRad~rad.

\claim{C-NOISE}The thermal noise is complex Gaussian. \claim{C-RNG}Random keys are derived with
a BLAKE2b hash of the seed, purpose and name, and the counter-based Philox4x32-10
generator~\cite{Salmon2011Philox} takes the absolute index (sample or step) as its counter,
so the compiled and reference backends produce identical integer streams (RT-14). Further
impairments are carrier-frequency offsets of TX and RX oscillators,
computed as phases of the absolute sample index (RT-09), Wiener phase noise with keyed
increments, and a front end with full-scale mapping, gain and optional clipping.

\subsection{What the release does not model}
% Limitations table (DESIGN 1.5 non-goals, plus the RisEnv backend limit of the committed
% code: sremu/ml/env.py accepts backend="batch" and raises for the IQ backends).
% Limitations table (DESIGN 1.5 non-goals, plus the RisEnv backend limit of the committed
% code: sremu/ml/env.py accepts backend="batch" and raises for the IQ backends).
\begin{table}[t]
\caption{What the current release of \sremu{} does not model or provide, and what it offers
instead.}
\label{tab:limitations}
\centering
\footnotesize
\setlength{\tabcolsep}{3pt}
\renewcommand{\arraystretch}{1.1}
\begin{tabular}{@{}>{\raggedright\arraybackslash}p{1.52in}>{\raggedright\arraybackslash}p{1.86in}@{}}
\toprule
Not modelled or not provided & Instead \\
\midrule
Unmodified UHD applications; a virtual USRP driver &
  UHD-shaped API and sockets; virtual-USRP emulation is the approach of
  ACHEM~\cite{Gurses2026ACHEM} \\
Hard real time; paced timing on Windows &
  Soft real time with counted misses; timing results from Linux hosts \\
Ray tracing, buildings, terrain profiles, several knife edges on one path, walls of order
  two, RIS--wall paths &
  Explicit scenes: LOS, flat ground, order-one walls, one knife edge per path, additive
  diffuse term \\
RIS--RIS paths; polarization & Single-bounce RIS paths; scalar co-polar fields \\
Attenuation of RIS hops by occluders; fading or multipath on RIS hops &
  Occluders act on the LOS path and its ground and wall images; RIS hops are deterministic
  free-space paths (optional flat-ground images) \\
Mutual coupling between differing neighbour states &
  Per-state tables, which hold coupling for identical neighbours (locally
  periodic)~\cite{ETSI_GR_RIS003_v121}; a coupling-hook protocol, implemented as no
  coupling \\
Angle-dependent tables outside one plane of incidence &
  Single-plane tables; off-plane components projected out, or refused on request \\
Switching, transient or timing behaviour compared with measurements &
  Transition profiles as an engineering choice; command records with requested and applied
  samples \\
Frequency-dependent $\Gamma$ filters; non-reciprocal tables &
  Signals wider than a table's validity bandwidth and non-reciprocal tables are rejected \\
Sampling-frequency offset, IQ imbalance, amplifier and converter models &
  CFO, phase noise, full-scale clipping, \texttt{sc16} saturation \\
Rate conversion between overlapping bands; several channels per stream &
  Linked streams share one rate; one channel per stream \\
Web server, dashboard, authentication &
  Static HTML report; loopback sockets by default \\
E2SM-RIS adapter; conformance to a standard; GPU IQ path; vectorized environments &
  Future work (Section~\ref{sec:future}) \\
Learning-based controller research; IQ backends of \texttt{RisEnv} &
  Controller protocol; \texttt{RisEnv} on the batch path; a UCB1 baseline \\
\bottomrule
\end{tabular}
\end{table}

\claim{C-LIMITS}Table~\ref{tab:limitations} lists what the current release does not model or
provide, and what it offers instead; Section~\ref{sec:disc-limits} names the four limits that
matter most for reading the results of this paper.
